\chapter{Επίλογος}

Το βασικό συμπέρασμα είναι ότι η κεντρική εκπαίδευση παραμένει σαφώς ισχυρότερη όταν η συγκέντρωση των δεδομένων είναι εφικτή. Η ομοσπονδιακή εκπαίδευση, ωστόσο, παρέχει μια πρακτική εναλλακτική όταν τα δεδομένα πρέπει να παραμείνουν κατανεμημένα, όπως συμβαίνει συχνά σε περιβάλλοντα δικτύων, παρόχων και συστημάτων όπου η ιδιωτικότητα, η ιδιοκτησία των δεδομένων ή το κόστος μεταφοράς αποτελούν περιοριστικούς παράγοντες.

% Η μεταβολή του πλήθους των \en{clients} δεν οδήγησε σε μονοτονική βελτίωση της απόδοσης. Η καλύτερη τιμή \en{F1-score} αυξήθηκε ελαφρά από \(0.8201\) στους 2 \en{clients} σε \(0.8245\) στους 4 \en{clients}, αλλά μειώθηκε σε \(0.8222\) στους 8 \en{clients}. Επομένως, η αύξηση του πλήθους των \en{clients} δεν αρκεί από μόνη της για καλύτερη απόδοση. Αντίθετα, αυξάνει τον κατακερματισμό των τοπικών δεδομένων και μπορεί να κάνει τη συνάθροιση των ενημερώσεων πιο δύσκολη.

Η επίδραση της στρατηγικής συνάθροισης είχε διαφορετικά αποτελέσματα ανάλογα με το πλήθος των \en{clients}. Στους 2 \en{clients}, η καλύτερη επίδοση προήλθε από τη \en{FedAvg} με \(\alpha=100\). Στους 4 και 8 \en{clients}, η καλύτερη επίδοση προήλθε από τη \en{FedAdam} με \(\alpha=100\). Αυτό δείχνει ότι η προσαρμοστική ενημέρωση του διακομιστή γίνεται πιο χρήσιμη όταν αυξάνεται το πλήθος των τοπικών ενημερώσεων και η διαδικασία συνάθροισης γίνεται πιο σύνθετη.

Η παράμετρος \(\alpha\) επηρέασε επίσης τη συμπεριφορά των πειραμάτων, αλλά όχι με απόλυτα ενιαίο τρόπο. Οι καλύτερες τιμές \en{F1-score} και στις τρεις ομάδες \en{clients} προέκυψαν με \(\alpha=100\), δηλαδή με κατανομή που προσεγγίζει περισσότερο την \en{IID} περίπτωση. Ωστόσο, ορισμένες ανομοιογενείς ρυθμίσεις με \(\alpha=0.8\) πέτυχαν ανταγωνιστικές τιμές, ιδιαίτερα ως προς το \en{recall}. Αυτό δείχνει ότι η μη ομοιόμορφη κατανομή δεν οδηγεί αναγκαστικά σε αποτυχία, αλλά μπορεί να μεταβάλει την ισορροπία ανάμεσα στον εντοπισμό ανωμαλιών και στην αποφυγή ψευδώς θετικών προβλέψεων.

Η εργασία επιβεβαιώνει επίσης ότι η προσαρμογή \en{LoRA} μπορεί να ενσωματωθεί σε ομοσπονδιακό πλαίσιο για ανίχνευση ανωμαλιών σε αρχεία καταγραφής. Το μοντέλο δεν χρειάζεται να εκπαιδεύσει όλες τις παραμέτρους του \en{BERT}, αλλά μόνο τις πρόσθετες παραμέτρους προσαρμογής. Αυτό καθιστά την προσέγγιση πιο κατάλληλη για κατανεμημένα περιβάλλοντα, στα οποία το κόστος επικοινωνίας και οι υπολογιστικοί πόροι των \en{clients} αποτελούν σημαντικούς περιορισμούς.

Από πρακτική άποψη, τα αποτελέσματα δείχνουν έναν σαφή συμβιβασμό. Η κεντρική εκπαίδευση προσφέρει την καλύτερη απόδοση, αλλά απαιτεί συγκέντρωση των δεδομένων. Η ομοσπονδιακή εκπαίδευση μειώνει αυτή την απαίτηση, αλλά συνοδεύεται από χαμηλότερη επίδοση. Για εφαρμογές όπου η μέγιστη δυνατή ακρίβεια είναι ο μοναδικός στόχος, η κεντρική προσέγγιση παραμένει προτιμότερη. Για εφαρμογές όπου η ιδιωτικότητα, η κατανεμημένη ιδιοκτησία των δεδομένων και η περιορισμένη μεταφορά πληροφορίας είναι κρίσιμοι παράγοντες, η ομοσπονδιακή προσέγγιση αποτελεί ρεαλιστική λύση.
